% Unit tests of RH: the problems with known answers and their derivations.
% Build: make -C docs/unit-tests  (pdflatex, twice for the references)

\documentclass[11pt,a4paper]{article}
\usepackage[T1]{fontenc}
\usepackage[utf8]{inputenc}
\usepackage{lmodern}
\usepackage[margin=25mm]{geometry}
\usepackage{amsmath,amssymb}
\usepackage{booktabs}
\usepackage{array}
\usepackage[round]{natbib}
\usepackage[hidelinks]{hyperref}
% journal macros used by ADS BibTeX exports
\providecommand{\jqsrt}{JQSRT}
\providecommand{\apj}{ApJ}
\providecommand{\aap}{A\&A}

\newcommand{\code}[1]{\texttt{#1}}
\newcommand{\dd}{\mathrm{d}}
\newcommand{\e}{\mathrm{e}}
\newcommand{\Kp}{K'}
\newcommand{\vI}{\boldsymbol{I}}
\newcommand{\vS}{\boldsymbol{S}}
\newcommand{\vP}{\boldsymbol{P}}
\newcommand{\eI}{\boldsymbol{e}_I}

\title{Unit tests of RH:\\ problems with known answers}
\author{RH modernization (\code{tests/unit})}
\date{\today}

\begin{document}
\maketitle

\begin{abstract}
The unit tests in \code{tests/unit} exercise single routines of RH on
problems whose answers are known, either in closed form or from an
independent high-precision computation. This document derives those
answers, states what each test checks and why the expected accuracy or
order of convergence is what it is. Where a test exposed a defect, the
defect is described with the evidence (identifiers such as B-24 refer to
\code{docs/MODERNIZATION.md}).
\end{abstract}

\tableofcontents

%======================================================================
\section{Approach}

RH has four tiers of tests:
\begin{description}
\item[unit] (\code{ctest -L unit}): single routines, known answers,
  milliseconds to seconds; run in CI, also with AddressSanitizer and
  UndefinedBehaviorSanitizer;
\item[end to end] (\code{ctest -L e2e}): whole runs of a program on
  problems with exact properties (section~\ref{sec:e2e});
\item[regression] (\code{tests/run\_tests.sh}): complete runs of the
  example problems compared with a baseline build. These runs stop after
  a few iterations and are \emph{not converged}; they detect change, not
  error;
\item[stress] (\code{tests/run\_stress.sh}): runs to convergence on
  harder atmospheres, for convergence behaviour, timing and memory.
\end{description}
Only the unit and end-to-end tiers can say whether an answer is \emph{right}. Their tests
follow a few rules:
\begin{itemize}
\item Reference values come from closed forms or from \code{mpmath} at
  50--60 digits, generated by the \code{make\_*\_reference.py} scripts
  into \code{*\_reference.h}, so that the C tests need no Python.
\item Where a method is exact for a class of inputs (a linear source
  function, polynomials up to some degree), the test requires exactness
  to rounding. This is the sharpest check there is: the first serious
  defect found this way (B-24, long characteristics one section short)
  showed up as a 7\% error where the answer must be exact.
\item Otherwise the test measures the order of convergence under grid
  refinement and the error at one resolution. Bounds are about twice the
  values measured with gcc on x86-64, and the measured values are written
  next to the bounds in the source.
\item Every test can print what it measures (\code{-v} or \code{--table}).
\end{itemize}

%======================================================================
\section[The Voigt and Faraday-Voigt functions]{The Voigt and Faraday--Voigt functions}
\label{sec:voigt}

\subsection{Definitions}
With damping parameter $a\ge 0$ and frequency $v$, both in Doppler
widths,
\begin{equation}
H(a,v) = \frac{a}{\pi}\int_{-\infty}^{\infty}
  \frac{\e^{-y^2}\,\dd y}{(v-y)^2+a^2},\qquad
F(a,v) = \frac{1}{\pi}\int_{-\infty}^{\infty}
  \frac{(v-y)\,\e^{-y^2}\,\dd y}{(v-y)^2+a^2},
\end{equation}
and $H + iF = w(z)$, $z = v + ia$, with the Faddeeva function
$w(z) = \e^{-z^2}\operatorname{erfc}(-iz)$. RH does not normalize:
$H(0,v) = \e^{-v^2}$ and $\int H\,\dd v = \sqrt\pi$. $H$ is even and $F$
odd in $v$.

\subsection{The implementation}
\code{voigt.c} combines (after \citealt{2018JQSRT.213...13S}):
\begin{itemize}
\item for $|v|+a\ge 8$, the Laplace continued fraction
\begin{equation}
w(z) = \frac{i}{\sqrt\pi}\,\cfrac{1}{z-\cfrac{1/2}{z-\cfrac{1}{z-\cfrac{3/2}{z-\cdots}}}}
\end{equation}
  to six levels, evaluated as the ratio $A_6/B_6$ of its convergents,
  $A_k = zA_{k-1}-\tfrac{k}{2}A_{k-2}$ and the same for $B_k$, so that
  there is a single complex division;
\item inside that region, the rational approximation of \citet{10.1137/0731077} with
  $N=32$: with $L = 2^{-1/4}N^{1/2}$, $r = 1/(L-iz)$ and
  $Z = (L+iz)\,r$,
  \begin{equation}
  w(z) \approx \Bigl(2\,p(Z)\,r + \tfrac{1}{\sqrt\pi}\Bigr)\,r,\qquad
  p(Z) = \sum_{n=0}^{N-1} a_n Z^n,
  \end{equation}
  the coefficients $a_n$ being the Fourier coefficients of the integrand
  after the substitution $y = L\tan(\theta/2)$
  (\code{make\_weideman\_coefficients.py}). $p$ is evaluated as two
  Horner chains in $Z^2$;
\item for $a<10^{-6}$, the expansion to first order in $a$. At $a=0$,
  $w(v) = \e^{-v^2} + iF_0(v)$ with $F_0 = (2/\sqrt\pi)\,D(v)$ ($D$ is
  Dawson's integral), and $w$ obeys $w'(z) = -2zw + 2i/\sqrt\pi$, so
  \begin{equation}
  w(v+ia) = w(v) + ia\,w'(v) + O(a^2)
  \;\Rightarrow\;
  \begin{cases}
  H = \e^{-v^2} + a\,(2vF_0 - 2/\sqrt\pi),\\
  F = F_0 - 2av\,\e^{-v^2},
  \end{cases}
  \end{equation}
  with a relative error $\sim 2a^2v^2$. Without this branch the absolute
  error of Weideman's approximation, $\sim10^{-14}$, would exceed $H$
  itself in the near wing of a weakly damped line, where $H\approx
  2a/(\sqrt\pi\,v^2)$ is tiny;
\item for $|z|>10^6$, the first term $i/(\sqrt\pi z)$, scaled so that
  $|z|^2$ cannot overflow (the convergents would).
\end{itemize}

\subsection{What \code{test\_voigt} checks}
Against $w(z)$ from \code{mpmath} on $a\in\{0,10^{-12},10^{-10},
10^{-8},10^{-6},\ldots,100\}$ and $0\le v\le 10^4$:
\begin{itemize}
\item the relative error of $H$ wherever $H\ge 10^{-12}H(a,0)$ (a strong
  line's wing matters against the continuum down to about there):
  measured $1.6\times10^{-6}$, largest for $a\sim10^{-6}$ just outside
  the Doppler core, where the first-order expansion and \citeauthor{10.1137/0731077}'s
  approximation meet;
\item the error relative to the line centre, $|H-H_\text{ref}|/H(a,0)$:
  $2.7\times10^{-10}$;
\item $F$ to within $10^{-9}$ of $\max|F|$ for the given $a$;
\item exact symmetry ($H$ even, $F$ odd), finite and positive $H$,
  continuity across the switches at $|v|+a=8$ and $a=10^{-6}$, finite
  values for $|z|$ up to $10^{200}$.
\end{itemize}
Underflowing reference values are flushed to zero, and every value is
checked with \code{isfinite} first: \code{MAX()} would hide a NaN. That
check found B-20 (NaN at $a=v=0$ in a former algorithm).

%======================================================================
\section{Numerical kernels}

\subsection{Linear systems (\code{test\_linalg})}
\code{ludcmp.c} (LU decomposition with implicit scaling and partial
pivoting) is tested on
\begin{itemize}
\item $A\,x = b$ with random, diagonally dominant $A$ ($N=1\ldots60$)
  and $b = Ax_\text{known}$: well conditioned, so $x$ must be accurate
  to a few times the machine precision;
\item a matrix with zeros on the diagonal, which needs pivoting, with
  integer entries: the solution and the determinant (the sign $d$ times
  the product of the diagonal of $LU$) are exact;
\item the Hilbert matrix scaled to integers,
  $A_{ij} = \operatorname{lcm}(1,\ldots,2N-1)/(i+j+1)$ ($0$-based), with
  condition number $1.5\times10^{10}$ at $N=8$: the error is then of
  order $\kappa\,\epsilon$, with or without iterative refinement
  (\code{improve});
\item the inverse from \code{LUbacksubst} on the columns of the unit
  matrix: $AA^{-1}=1$.
\end{itemize}

\subsection{Interpolation (\code{test\_interp})}
\paragraph{Table look-up.} \code{Hunt} and \code{Locate} must return
$i$ with $x_i\le x<x_{i+1}$ for values at, between, below and above the
nodes, ascending and descending tables, tables of $1,2,3$ and many
elements, and every starting guess of \code{Hunt}. (This found B-23:
\code{Hunt} one index low when hunting downward onto a node.)

\paragraph{Linear interpolation.} Exact for linear data; for smooth $f$
on a spacing $h$,
\begin{equation}
f(x) - L(x) = \tfrac12 f''(\xi)\,(x-x_0)(x-x_1),\qquad
|f-L| \le \tfrac18 h^2\max|f''|,
\end{equation}
since $|(x-x_0)(x-x_1)|\le h^2/4$: second order. \code{BiLinear} is
exact for bilinear data. (B-21: \code{BiLinear} reads past its table at
the upper edge.)

\paragraph{Cubic splines.} The natural cubic spline has $S''=0$ at both
ends. For $\sin x$ on $[0,\pi]$, where $f''=0$ at the ends, the boundary
conditions are exact and the spline converges at fourth order
everywhere. For $\e^x$ they are wrong by $O(1)$ in $S''$, which costs
$O(h^2)$ in $S$ near the ends, decaying into the interior: fourth order
in the interior, second overall. The test checks both, and the spline
against the solution of its defining conditions computed with
\code{mpmath}.

\paragraph{Exponential splines.} Between nodes the exponential spline
with tension $\sigma$ satisfies $y'''' - \sigma^2y'' = 0$, with basis
$\{1,x,\e^{\sigma x},\e^{-\sigma x}\}$. It tends to the cubic spline for
$\sigma\to0$ and to linear interpolation for $\sigma\to\infty$; the test
checks both limits and two tensions against \code{mpmath}.

\paragraph{Cubic convolution.} The kernel of \citet{1981ITASS..29.1153K}
\begin{equation}
u(s) = \begin{cases}
\tfrac32|s|^3-\tfrac52|s|^2+1, & |s|\le1,\\
-\tfrac12|s|^3+\tfrac52|s|^2-4|s|+2, & 1<|s|<2,\\
0 & \text{otherwise}
\end{cases}
\end{equation}
is a partition of unity and reproduces quadratics, so interpolation is
third order; with \citeauthor{1981ITASS..29.1153K}' extrapolation of the data at the edges
($f_{-1} = 3f_0-3f_1+f_2$) it is exact for biquadratic data in every
cell, including the boundary cells. All three properties are checked.

\subsection{Quadrature (\code{test\_quadrature})}
\paragraph{Gauss--Legendre.} The $n$-point rule integrates polynomials
of degree $\le 2n-1$ exactly. For $x^{2n}$ on $[-1,1]$ the remainder
\begin{equation}
E_n[f] = \frac{2^{2n+1}(n!)^4}{(2n+1)\,[(2n)!]^3}\,f^{(2n)}(\xi)
\end{equation}
is known exactly, $f^{(2n)} = (2n)!$, so the test checks the error
$2^{2n+1}(n!)^4/\bigl((2n+1)[(2n)!]^2\bigr)$ as well as exactness below.
Also: nodes inside the interval, increasing and symmetric; positive
weights summing to the length; the closed forms for $n=2,3,4$
($\pm1/\sqrt3$; $0,\pm\sqrt{3/5}$ with weights $8/9, 5/9$;
$\pm\sqrt{3/7\mp\frac27\sqrt{6/5}}$ with weights
$(18\pm\sqrt{30})/36$); a reversed interval changes the sign.

\paragraph{Weights of piecewise integration.} \code{w3.c} returns
\begin{equation}
w_k(\Delta\tau) = \int_0^{\Delta\tau} t^k\e^{-t}\,\dd t = \gamma(k+1,\Delta\tau),\qquad k=0,1,2,
\end{equation}
the lower incomplete gamma function, used by the linear short
characteristics. The closed forms
($w_0 = 1-\e^{-\Delta\tau}$, $w_1 = w_0 - \Delta\tau\,\e^{-\Delta\tau}$,
$w_2 = 2w_1 - \Delta\tau^2\e^{-\Delta\tau}$) cancel for small
$\Delta\tau$; below $\Delta\tau = 1$ the code sums the series
\begin{equation}
\gamma(m,t) = t^m\e^{-t}\sum_{k\ge0}\frac{t^k}{m(m+1)\cdots(m+k)},
\end{equation}
whose terms are all positive. The test compares with \code{mpmath} on
$0\le\Delta\tau\le1000$, across all switches. (N-18: the former series
was wrong by up to $2\times10^{-6}$ near $\Delta\tau=5\times10^{-4}$.)

\subsection{Special functions (\code{test\_functions})}
\paragraph{Exponential integrals.} $E_1$ uses formula 5.1.53 of \citet{1972hmfw.book.....A} (polynomial, absolute error $<2\times10^{-7}$ for $x\le1$) and
5.1.56 (rational, $x\e^xE_1$ to $2\times10^{-8}$ relative for $x\ge1$);
the test holds them to those accuracies against \code{mpmath}.
For $x<1$, $E_2 = \e^{-x}-xE_1$ (5.1.14). At larger $x$ this cancels,
since $\e^{-x}$ and $xE_1\approx\e^{-x}(1-1/x)$ nearly agree, and the
relative error of $E_2$ grows as $x$ times that of $E_1$; for $x\ge1$
$E_2$ therefore comes from the continued fraction
\begin{equation}
E_n(x) = \e^{-x}\left(\frac{1}{x+n-}\,\frac{1\cdot n}{x+n+2-}\,
\frac{2(n+1)}{x+n+4-}\cdots\right)
\end{equation}
(5.1.22), evaluated by the modified Lentz method, which the test holds
to $10^{-14}$. Both functions are zero only where $\e^{-x}/x$ underflows ($x\approx740$)
(N-22; before, both returned $0$ above $x=80$).

\paragraph{Planck function.} Against the formula in high precision with
the constants of \code{constant.h} (the denominator is evaluated with
\code{expm1}, so the relative error does not grow as $1/x$ at
radio wavelengths); the Wien ($h\nu\gg kT$:
$B_\nu\to 2h\nu^3c^{-2}\e^{-h\nu/kT}$) and Rayleigh--Jeans
($B_\nu\to2\nu^2kT/c^2$) limits; and
\begin{equation}
\int_0^\infty B_\nu\,\dd\nu = \frac{2k^4T^4}{h^3c^2}\int_0^\infty\frac{x^3\,\dd x}{\e^x-1}
= \frac{2k^4T^4}{h^3c^2}\,\frac{\pi^4}{15} = \frac{\sigma T^4}{\pi},
\end{equation}
by Gauss--Legendre quadrature in $x=h\nu/kT$.

\paragraph{3$j$ and 6$j$ symbols.} \code{w3js} against
$\begin{pmatrix}j&j&0\\m&-m&0\end{pmatrix} = (-1)^{j-m}/\sqrt{2j+1}$;
$\begin{pmatrix}j_1&j_2&j_3\\0&0&0\end{pmatrix}$, zero for odd
$J=j_1+j_2+j_3$ and otherwise
\begin{equation}
(-1)^{J/2}\sqrt{\frac{(J-2j_1)!\,(J-2j_2)!\,(J-2j_3)!}{(J+1)!}}\;
\frac{(J/2)!}{(J/2-j_1)!\,(J/2-j_2)!\,(J/2-j_3)!};
\end{equation}
the stretched case $j_3=j_1+j_2$; the orthogonality
$\sum_{m_1m_2}\begin{pmatrix}j_1&j_2&j_3\\m_1&m_2&m_3\end{pmatrix}
\begin{pmatrix}j_1&j_2&j_3'\\m_1&m_2&m_3\end{pmatrix} =
\delta_{j_3j_3'}/(2j_3+1)$; the symmetries under column permutations and
$m\to-m$; the selection rules. \code{w6js} against
$\{a\,b\,c;0\,c\,b\} = (-1)^{a+b+c}/\sqrt{(2b+1)(2c+1)}$, the sum rules
\begin{gather}
\sum_x(2x+1)\begin{Bmatrix}a&b&x\\a&b&f\end{Bmatrix} = (-1)^{2(a+b)},\\
\sum_x(-1)^{a+b+x}(2x+1)\begin{Bmatrix}a&b&x\\b&a&f\end{Bmatrix} =
\sqrt{(2a+1)(2b+1)}\,\delta_{f0},
\end{gather}
and the orthogonality of the 6$j$ symbols. (B-22: \code{w3js} accepts
invalid arguments and overflows from $j=50$.)

%======================================================================
\section{Formal solution in one dimension}
\label{sec:formal1d}

\subsection{Slabs with a polynomial source function}
In a plane-parallel atmosphere with vertical optical depth $\tau$
(increasing inward) and direction cosine $\mu>0$, the outgoing
($I^+$) and incoming ($I^-$) intensities obey
\begin{equation}
\mu\frac{\dd I^+}{\dd\tau} = I^+ - S,\qquad
-\mu\frac{\dd I^-}{\dd\tau} = I^- - S,
\end{equation}
with formal solutions
\begin{align}
I^+(\tau) &= \int_\tau^\infty S(t)\,\e^{-(t-\tau)/\mu}\,\frac{\dd t}{\mu}, \\
I^-(\tau) &= \int_{\tau_0}^\tau S(t)\,\e^{-(\tau-t)/\mu}\,\frac{\dd t}{\mu}
  \qquad (I^-(\tau_0)=0 \text{ at the top } \tau_0).
\end{align}
For a polynomial $S(\tau) = \sum_n s_n\tau^n$, repeated integration by
parts terminates:
\begin{align}
I^+(\tau) &= \sum_{j\ge0}\mu^j S^{(j)}(\tau), \label{eq:Iplus}\\
I^-(\tau) &= P^-(\tau) - \e^{-(\tau-\tau_0)/\mu}P^-(\tau_0),\qquad
P^-(\tau) = \sum_{j\ge0}(-\mu)^jS^{(j)}(\tau), \label{eq:Iminus}
\end{align}
with $S^{(j)}(\tau) = \sum_n s_n\frac{n!}{(n-j)!}\tau^{n-j}$.
Eq.~(\ref{eq:Iplus}) is the semi-infinite solution; a finite slab
reproduces it when its lower boundary is irradiated with $I^+(\tau_\text{max})$.

Just below the top, $I^-\to0$ is the small difference of two $O(1)$
terms. \code{slab\_exact.h} therefore evaluates
\begin{equation}
I^-(\tau) = \bigl[P^-(\tau)-P^-(\tau_0)\bigr]
  - \operatorname{expm1}\bigl(-(\tau-\tau_0)/\mu\bigr)\,P^-(\tau_0),
\end{equation}
with $\tau^p-\tau_0^p = (\tau-\tau_0)\sum_{i<p}\tau^{p-1-i}\tau_0^i$, so
that no cancellation remains. (Without it the reference itself lost
digits near the top.)

The test atmospheres have $N$ points logarithmic in $\tau$ from
$10^{-4}$ to $10^2$. The opacity is either constant ($\chi=1$, height
$z=-\tau$) or exponential in height ($\chi=\tau=\e^{-z}$), so that the
solvers must integrate $\chi$ over height to obtain $\tau$.

\subsection{Cubic Bezier integration}
\label{sec:bezier}
Over an interval of optical depth $\delta$ from the upwind point $u$ to
the central point $c$, with $x\in[0,1]$,
\begin{equation}
I_c = \e^{-\delta}I_u + \delta\int_0^1 S(x)\,\e^{-\delta(1-x)}\,\dd x .
\end{equation}
With $S$ a cubic Bezier curve,
$S(x) = S_u B_0 + C_u B_1 + C_c B_2 + S_c B_3$ and the Bernstein
polynomials $B_0 = (1-x)^3$, $B_1 = 3x(1-x)^2$, $B_2 = 3x^2(1-x)$,
$B_3 = x^3$,
\begin{equation}
I_c = \varepsilon I_u + \alpha S_u + \beta S_c + \gamma C_u + \theta C_c,
\end{equation}
\begin{equation}
\begin{aligned}
\varepsilon &= \e^{-\delta}, &
\alpha &= \frac{6-(6+6\delta+3\delta^2+\delta^3)\varepsilon}{\delta^3}, &
\beta  &= \frac{\delta^3-3\delta^2+6\delta-6+6\varepsilon}{\delta^3},\\
&& \gamma &= \frac{3\,[2\delta-6+(6+4\delta+\delta^2)\varepsilon]}{\delta^3}, &
\theta &= \frac{3\,[6-4\delta+\delta^2-(6+2\delta)\varepsilon]}{\delta^3}
\end{aligned}
\end{equation}
\citep{2013ApJ...764...33D}. $\gamma$ belongs to the
upwind control point and $\theta$ to the central one: for large
$\delta$, $\theta\approx3/\delta$ (near the central point, little
attenuated) and $\gamma\approx6/\delta^2$. The control points of the
cubic Hermite interpolant with derivatives $S'_u$, $S'_c$ are
\begin{equation}
C_u = S_u + \tfrac{\delta}{3}S'_u,\qquad C_c = S_c - \tfrac{\delta}{3}S'_c .
\end{equation}
RH estimates $S'_c$ from the upwind and downwind neighbours with
the monotonic limiter of \citet{1990A&A...239..443S}: with the one-sided slopes
$s_u = (S_c-S_u)/\delta_u$, $s_d = (S_d-S_c)/\delta_d$ and the parabolic
estimate $p = (s_u\delta_d+s_d\delta_u)/(\delta_u+\delta_d)$,
\begin{equation}
S'_c = \begin{cases}0 & s_us_d\le0,\\
\operatorname{sign}(p)\,\min(2|s_u|,2|s_d|,|p|) & \text{otherwise},
\end{cases}
\end{equation}
and carries it to the next interval as the upwind derivative.

\paragraph{Small $\delta$.} The closed forms cancel: each weight tends
to $\delta/4$ while its numerator is $O(\delta^4)$, losing about
$\log_{10}\delta^{-3}$ digits ($10^{-9}$ relative at $\delta=0.05$).
Below $|\delta| = 0.2$ \code{Bezier3\_coeffs} sums Taylor series with
2, 4, 7 or 9 terms (Estrin's scheme), and below $0.05$ obtains
$\varepsilon$ from the sum of the weights; above $\delta=40$ the closed
forms are expanded in $1/\delta$. \code{test\_formal\_1d} checks the
weights against \code{mpmath} at 60 digits for
$10^{-8}\le\delta\le10^3$ (\code{make\_bezier\_reference.py}). (N-17:
the former truncated series was wrong by up to $1.2\times10^{-5}$.)

\subsection{What \code{test\_formal\_1d} checks, and why}
\paragraph{Exactness.} If $S$ is linear in $\tau$, \citeauthor{1990A&A...239..443S}'s estimate
returns the exact slope (both one-sided slopes and $p$ are equal), the
Bezier curve reproduces $S$, and the linear solver's interpolant is
exact as well: both must give $I^\pm$ to rounding (measured
$10^{-15}$--$10^{-13}$). With exponential opacity this holds only for the
integration in $\tau$, not in height, so the exactness check uses
constant opacity.

\paragraph{Order of convergence.} On the cubic
$S = 1+\tau+\tfrac12\tau^2+\tfrac1{10}\tau^3$ and grids of $41$ to
$321$ points:
\begin{center}
\begin{tabular}{lcccc}
\toprule
& \multicolumn{2}{c}{order, emergent / interior} & \multicolumn{2}{c}{error at $N=161$, emergent}\\
method & constant $\chi$ & exponential $\chi$ & constant & exponential\\
\midrule
linear & 2.00 / 1.91 & 2.00 / 1.89 & $9.7\times10^{-4}$ & $1.4\times10^{-4}$\\
Bezier3 & 4.00 / 2.16 & 4.00 / 1.84 & $9.3\times10^{-7}$ & $3.4\times10^{-7}$\\
Feautrier & 2.00 & 1.99 & $4.4\times10^{-4}$ & $3.9\times10^{-4}$\\
Feautrier--Hermite & 3.97 / 4.00 & 1.98 & $3.6\times10^{-7}$ & $8.3\times10^{-4}$\\
\bottomrule
\end{tabular}
\end{center}
Why Bezier3 is fourth order for the emergent intensity although its
derivative estimate is only second order: the error of the parabolic
estimate, $\tfrac16\delta_u\delta_dS'''$, enters the interval
$[u,c]$ through $C_c$ with weight $-\tfrac\delta3\theta$ and the next
interval through $C_u$ with weight $+\tfrac\delta3\gamma$. For small
$\delta$, $\gamma$ and $\theta$ differ by $O(\delta^2)$, so the two
contributions cancel to higher order along the ray. At the two ends of
the ray the derivative is one-sided (the secant), which breaks the
cancellation; its error is carried inward and attenuated, which is why
the interior converges more slowly than the emergent intensity.
Feautrier--Hermite \citep{1976JQSRT..16..931A} is fourth order with constant opacity, but its
trapezoidal optical depth makes it second order when $\chi$ varies.

\paragraph{The diagonal operator.} MALI needs $\Psi_k = \partial
I_k/\partial S_k$ (per unit emissivity after division by $\chi_k$). The
linear and Feautrier solvers are linear in $S$, so $\Psi_k$ must equal
the intensity at $k$ for $S=\boldsymbol e_k$ and no incident radiation:
measured to $2\times10^{-16}$ and $1.9\times10^{-15}$. Bezier3 is not
linear in $S$ (the limiter); its $\Psi = \beta+\theta$ is compared with a
finite-difference $\partial I_k/\partial S_k$ and only reported:
it differs by up to 28\% (N-19, section~\ref{sec:psi}).

\subsection{The exact diagonal of the Bezier operator (N-19)}
\label{sec:psi}
$\beta+\theta$ is the derivative of $I_c$ with respect to $S_c$ at
\emph{fixed control points}. But $C_c = S_c-\tfrac\delta3S'_c$ and
$C_u$ of the same interval contains $S'_u$, which depends on $S_c$ (it
is the derivative at the previous point, whose downwind neighbour is
$c$); and $I_u$ depends on $S_c$ through the control point $C_c$ of the
previous interval. With $k-1\to k$ the interval ending at $k$,
\begin{equation}
\frac{\partial I_k}{\partial S_k} =
\beta_k + \theta_k\Bigl(1-\frac{\delta_k}{3}\frac{\partial S'_k}{\partial S_k}\Bigr)
+ \Bigl(\gamma_k\frac{\delta_k}{3} - \varepsilon_k\theta_{k-1}\frac{\delta_{k-1}}{3}\Bigr)
  \frac{\partial S'_{k-1}}{\partial S_k},
\end{equation}
where the partial derivatives of \citeauthor{1990A&A...239..443S}'s estimate follow from its
active branch:
\begin{equation}
\Bigl(\frac{\partial S'_c}{\partial S_u},\frac{\partial S'_c}{\partial S_c},
\frac{\partial S'_c}{\partial S_d}\Bigr) =
\begin{cases}
\Bigl(-\frac{\delta_d/\delta_u}{\delta_u+\delta_d},\;
\frac{\delta_d/\delta_u-\delta_u/\delta_d}{\delta_u+\delta_d},\;
\frac{\delta_u/\delta_d}{\delta_u+\delta_d}\Bigr) & S'_c = p,\\[1ex]
\bigl(-2/\delta_u,\;2/\delta_u,\;0\bigr) & S'_c = 2s_u,\\
\bigl(0,\;-2/\delta_d,\;2/\delta_d\bigr) & S'_c = 2s_d,\\
(0,0,0) & S'_c = 0.
\end{cases}
\end{equation}
This costs a few multiplications per point. A prototype agrees with the
finite-difference derivative to the accuracy of the latter
($3\times10^{-5}$, limited by roundoff for small steps and by the kinks
of the limiter for large ones).

%======================================================================
\section{Formal solution in two and three dimensions}

\subsection{The uniform slab}
In a horizontally uniform, horizontally periodic slab every column must
reproduce the one-dimensional solution along the ray's vertical
direction cosine $\mu_z$, eqs.~(\ref{eq:Iplus}) and (\ref{eq:Iminus}).
The tests use the 1-D atmosphere of section~\ref{sec:formal1d} as the
vertical grid ($6$ columns in 2-D, $4\times4$ in 3-D) and several rays
in different quadrants or octants, including a vertical one.

\subsection{Meshes}
A short characteristic through a grid point meets the layer above or
below at a horizontal distance $\Delta z\,|\mu_x/\mu_z|$. If that is less
than the column spacing $\Delta x$, the upwind point lies on a layer,
between two columns; otherwise it lies on a column line, between two
layers, and S, $\chi$ and $I$ are interpolated vertically (quadratic,
monotonic).
\begin{description}
\item[``layers''] $\Delta x = 10^8$: every upwind and downwind point lies
  on a layer, where interpolation of horizontally uniform data is exact;
  the multi-D solvers must reduce to the 1-D ones.
\item[``columns''] $\Delta x = 0.05$: deep layers (large $\Delta z$) are
  crossed at column lines, so vertical interpolation and long
  characteristics are exercised. In 2-D also ``uneven'', with columns
  alternately $0.03$ and $0.07$ wide.
\end{description}
Any of the interpolations reproduces a function linear in height, so
$I^+$ for linear $S$ (which is linear in $\tau$, eq.~\ref{eq:Iplus}) must
be exact on every mesh; $I^-$ is not linear near the top and is exact
only on ``layers''.

\subsection{Long characteristics (B-24)}
With periodic horizontal boundaries the first point of a row needs the
intensity at an upwind point in the \emph{same} row, which has not been
computed yet. RH then follows the ray back across $N_\text{lc}$ column
crossings to a point on a layer and integrates forward along this long
characteristic (\code{SolveLong}, \code{SolveLongStokes}). The loop
treated the last crossing as the local point and returned the intensity
one crossing upwind of the true upwind point, skipping the last section.
For linear $S$, where the answer must be exact, the emergent intensity
was 7\% (2-D) and 11\% (3-D) wrong and identical columns differed by
15\%. The fix makes the downwind end of the last section the local
point; $I^+$ for linear $S$ is now exact to $2\times10^{-15}$.

\subsection{Orders of convergence}
On ``layers'' the linear solver is second order, as in 1-D. Bezier3 was
also only second order (N-24): at the upwind point it used the secant
slope $(S_c-S_u)/\delta$ instead of a derivative, so the Bezier curve was
exact only for linear $S$ (emergent error $4.8\times10^{-4}$ at
$N_z=161$, order 1.99). It now uses the slope of the quadratic Bezier
curve through the two points with the local derivative $S'_c$,
\begin{equation}
S'_u = 2\,\frac{S_c-S_u}{\delta} - S'_c ,
\end{equation}
which is the cubic form of the quadratic Bezier interpolant with one
control point $S_c - \tfrac\delta2 S'_c$ and is exact for quadratic $S$;
the same for $\chi$ (and, polarized, $\vS$ and $K$). The Steffen limiter
keeps $S'_c$ between 0 and twice the secant, so $S'_u$ lies there too and
the curve stays monotonic. On ``layers'' the order is now 2.82, emergent
$2.5\times10^{-5}$ at $N_z=161$ (2-D and 3-D); the 1-D solver, which
carries a centred derivative to the upwind point, is fourth order. On
``columns'' the vertical interpolation of $I$ limits Bezier3 to order
2.6 in 2-D, and the long characteristics both solvers to order
$\approx1.5$ in 3-D at these resolutions. Columns of the uniform slab agree to $8\times10^{-3}$
on ``columns'' (2-D, the two layers nearest the bottom) and exactly on
``layers''.

%======================================================================
\section{Spherical symmetry}
\code{test\_formal\_sphere} solves a shell of thickness $H=1$ around a
core of radius $R=1$ with constant opacity $\chi=10$ and constant source
function $S = B_\nu(T)$. Along a ray with impact parameter $p$ the path
coordinate is $s = \sqrt{r^2-p^2}$, with $s_0$ at the outer radius.
\begin{itemize}
\item Shell rays ($p>R$) pass the tangent point $s=0$:
  \begin{equation}
  I^-(s) = S\bigl(1-\e^{-\chi(s_0-s)}\bigr),\qquad
  I^+(s) = S\bigl(1-\e^{-\chi(s_0+s)}\bigr),
  \end{equation}
  emergent $S(1-\e^{-2\chi s_0})$.
\item Core rays ($p<R$) end on a thermalized core, $I^+ = B = S$ there,
  and with $S=B$ everywhere $I^+(s) = S$; $I^-$ as above.
\end{itemize}
Feautrier returns $P=\tfrac12(I^++I^-)$ and the emergent $I^+$. The
grid along a ray is where it crosses the radial shells; near the tangent
point $s\approx\sqrt{2p\,(r-p)}$, so the sampling is coarse there and the
measured orders, 1.4--1.6 for standard Feautrier and 1.2--1.5 for
Feautrier--Hermite, stay below the formal 2 and 4. Rays whose tangent
point lies above radial optical depth $0.1$ (grazing rays near the top)
are not refined by refining the radial grid: their tangent points move
toward the surface and the outermost always has three points. Their
errors (up to $7\times10^{-3}$) do not converge and are bounded
separately.

%======================================================================
\section{Polarized formal solution}
\label{sec:stokes}

\subsection{Transfer equation and propagation matrix}
For the Stokes vector $\vI = (I,Q,U,V)^T$, in optical depth of Stokes $I$
along the ray,
\begin{equation}
\mu\frac{\dd\vI}{\dd\tau} = K(\vI-\vS),\qquad K = 1+\Kp,
\end{equation}
with the reduced propagation matrix of \code{StokesK()},
\begin{equation}
K = \begin{pmatrix}
1 & \eta_Q & \eta_U & \eta_V\\
\eta_Q & 1 & \rho_V & -\rho_U\\
\eta_U & -\rho_V & 1 & \rho_Q\\
\eta_V & \rho_U & -\rho_Q & 1
\end{pmatrix},
\end{equation}
all elements divided by $\chi_I$. In LTE, and in the tests, the source
vector is unpolarized, $\vS = (S,0,0,0)^T$; RH is given the emission
vector over $\chi_I$, which is $K\vS$.

\subsection[Exact solution for constant K]{Exact solution for constant $K$}
For constant $K$ and polynomial $S$ the scalar derivation carries over
with $\mu\to\mu K^{-1}$ (the series below terminate, and $K$ commutes
with its own powers):
\begin{align}
\vI^+(\tau) &= \sum_{j\ge0}(\mu K^{-1})^j S^{(j)}(\tau)\,\eI, \\
\vI^-(\tau) &= \vP^-(\tau) - \e^{-K(\tau-\tau_0)/\mu}\,\vP^-(\tau_0),\qquad
\vP^-(\tau) = \sum_{j\ge0}(-\mu K^{-1})^j S^{(j)}(\tau)\,\eI .
\end{align}
Substitution confirms it: $\mu\,\dd\vI^+/\dd\tau =
\sum_j\mu^{j+1}K^{-j}S^{(j+1)}\eI = K(\vI^+-S\eI)$. As in 1-D,
$\vI^-$ is evaluated without cancellation, with
$\vP^-(\tau)-\vP^-(\tau_0)$ from the factored polynomial differences and
$\e^X-1$ by its Taylor series after scaling $X$ by $2^{-s}$, followed by
$s$ squarings $E\to E(E+2)$ (\code{stokes\_exact.h}). The test values
($\eta_{Q,U,V} = 0.12, -0.05, 0.25$, $\rho_{Q,U,V} = 0.04, 0.09, 0.3$)
are typical of a strong Zeeman-split line core; $K$ is positive definite
so the bottom boundary condition of the solvers ($Q=U=V=0$, unlike the
semi-infinite solution) is forgotten within a few optical depths. The
tests compare only $\tau\le10$.

\subsection{DELO-Bezier}
The polarized solvers \citep[DELO;][]{1989ApJ...339.1093R} write the transfer equation along the ray as
\begin{equation}
\frac{\dd\vI}{\dd t} = \vS_\text{eff} - \vI,\qquad
\vS_\text{eff} = \vS_K - \Kp\vI,
\end{equation}
with $\vS_K = K\vS$ the emission vector and $t$ the optical depth along
the propagation, and integrate $\vS_\text{eff}$ with the cubic Bezier
weights of section~\ref{sec:bezier}, treating $\vI_c$ implicitly. With
$\dd\Kp/\dd t$ written $\dot K$,
\begin{equation}
\vS_\text{eff}' = \vS_K' - \dot K\vI - \Kp\vS_K + (\Kp{}^2+\Kp)\vI,
\end{equation}
so the control points
$\boldsymbol C_u = \vS_{\text{eff},u}+\tfrac\delta3\vS'_{\text{eff},u}$
(weight $\gamma$) and
$\boldsymbol C_c = \vS_{\text{eff},c}-\tfrac\delta3\vS'_{\text{eff},c}$
(weight $\theta$) give
\begin{multline}
\Bigl[1+\beta\Kp_c+\theta\Bigl(\Kp_c+\tfrac\delta3\bigl(\Kp_c{}^2-\dot K_c+\Kp_c\bigr)\Bigr)\Bigr]\vI_c
= \Bigl[\varepsilon-\alpha\Kp_u+\gamma\Bigl(\tfrac\delta3\bigl(\Kp_u{}^2-\dot K_u+\Kp_u\bigr)-\Kp_u\Bigr)\Bigr]\vI_u\\
+ \Bigl[\alpha+\gamma\bigl(1-\tfrac\delta3\Kp_u\bigr)\Bigr]\vS_{K,u}
+ \Bigl[\beta+\theta\bigl(1+\tfrac\delta3\Kp_c\bigr)\Bigr]\vS_{K,c}
+ \tfrac\delta3\bigl(\gamma\vS'_{K,u}-\theta\vS'_{K,c}\bigr).
\end{multline}

\subsection{What the polarized tests check}
\begin{itemize}
\item With $\Kp=0$ the polarized solver must equal the unpolarized
  Bezier3 solver and give $Q=U=V=0$: measured $1.6\times10^{-15}$ (1-D)
  and $4.6\times10^{-16}$ (2-D, 3-D).
\item For linear $S$ and constant $K$, $\vI^+ = (S + \mu K^{-1}S')\eI$ is
  linear in $\tau$ and the solvers are exact: $7\times10^{-14}$ (1-D),
  $3\times10^{-13}$ (2-D, 3-D, all meshes).
\item For the cubic $S$ and constant $K$ (1-D): the emergent Stokes
  vector converges at order 4.00, $9.9\times10^{-7}$ at $N=161$ (the
  unpolarized solver: $9.3\times10^{-7}$); $\vI^-$ in the interior at
  order 2.9 (one-sided derivatives at the ends, as unpolarized).
\item For depth-dependent $K$, $\Kp(\tau) = \Kp_0/(1+\tau/2)$, there is
  no closed form; the order is measured from successive grids,
  $\log_2\bigl(|\vI_N-\vI_{N/2}|/|\vI_{2N}-\vI_N|\bigr)$: 4.00.
\item In 2-D, on ``layers'', the emergent Stokes vector for the cubic $S$
  converges at order 2.80, $3.3\times10^{-5}$ at $N_z=161$, like the
  unpolarized solver (N-24; second order before).
\end{itemize}

\paragraph{B-25.} These checks found two errors in the DELO-Bezier
solvers of all three geometries. (i) The derivative terms of the control
points were paired with the wrong weights,
$\tfrac\delta3(\gamma\vS'_{K,c}-\theta\vS'_{K,u})$ instead of
$\tfrac\delta3(\gamma\vS'_{K,u}-\theta\vS'_{K,c})$: with $\Kp=0$ the
polarized solver differed from the unpolarized one by 2.4\% at $N=81$,
and was second order. (ii) The sign of $\dot K_c$ on the left-hand side
was $+$: with depth-dependent $K$ the solver was first order (1.02).
Neither shows for linear $S$ and constant $K$, the case where the
derivatives are all equal or zero.

\paragraph{B-26.} In the last layer of a ray (the top layer for the
outgoing direction, which gives the emergent spectrum) the 2-D and 3-D
polarized solvers used the source vector $\vS_c$ before setting it: in
2-D the value of the last point of the layer below, in 3-D an
uninitialized local variable. With $\Kp=0$ the bottom layer of $\vI^-$
differed from the unpolarized solver by 38\%; for linear $S$ the
emergent vector was off by $2\times10^{-9}$ (the top interval is
optically very thin in the test).

%======================================================================
\section{Zeeman effect: splitting, profiles and Milne--Eddington}
\label{sec:zeeman}

\code{test\_zeeman} checks how RH turns an atomic line and a magnetic
field into the propagation matrix: the Zeeman pattern (\code{zeeman.c}),
the absorption and magneto-optical profiles (\code{profile.c}), the field
angles (\code{Bproject}), and, through a Milne--Eddington slab, the
emergent Stokes vector. The expectations are derived here from first
principles, following the conventions of \citet{2004ASSL..307.....L},
and only then compared with the code.

\subsection{Zeeman pattern}
In a field $B$ a level with total angular momentum $J$ and Land\'e factor
$g$ splits into sublevels $M=-J\ldots J$ with energies
$E_0+g\mu_\mathrm{B}BM$. A component $(M_u,M_l)$ of a line, with
$q = M_l-M_u\in\{-1,0,1\}$, absorbs at
$\nu = \nu_0+\nu_L(g_uM_u-g_lM_l)$, $\nu_L = eB/(4\pi m_ec)$. In the
reduced wavelength of RH, $v = (\lambda-\lambda_0)/\Delta\lambda_D$, its
centre is at
\begin{equation}
v_{M_uM_l} = (g_lM_l-g_uM_u)\,v_B,\qquad
v_B = \frac{\Delta\lambda_B}{\Delta\lambda_D},\quad
\Delta\lambda_B = \frac{\lambda_0^2\,eB}{4\pi m_ec^2},
\end{equation}
so $q=-1$ components are blue-shifted for $g>0$. The relative strengths
are
\begin{equation}
S_q(M_u,M_l) = 3\begin{pmatrix}J_u&J_l&1\\-M_u&M_l&-q\end{pmatrix}^2,
\qquad \sum_{M_uM_l}S_q = 1 \text{ for each } q,
\end{equation}
and in LS coupling $g = 1+[J(J+1)+S(S+1)-L(L+1)]/[2J(J+1)]$. The
effective Land\'e factor
\begin{equation}
\bar g = \tfrac12(g_u+g_l)+\tfrac14(g_u-g_l)\,[J_u(J_u+1)-J_l(J_l+1)]
\end{equation}
is the barycentre of the $q=+1$ components,
$\bar g = \sum S_{+1}(g_lM_l-g_uM_u)$; the test checks both forms.

\subsection{The propagation matrix from the normal modes}
With $H_q$ and $F_q$ the strength-weighted sums of the Voigt and
Faraday--Voigt functions over the components of each $q$,
$H_q(v) = \sum S_q H(a, v-v_{M_uM_l})$, the medium's response follows
from two facts.

\paragraph{Which components absorb which polarization.} A photon of
positive helicity carries angular momentum $+\hbar$ along its direction
of propagation; absorbed along $\boldsymbol B$, it raises $M$ by one:
$M_u=M_l+1$, $q=-1$. In the IAU convention positive Stokes $V$ is
positive helicity (the electric vector turns counterclockwise as seen by
the observer). Classically the same follows from Larmor's theorem: an
electron circling counterclockwise about $\boldsymbol B$ (seen from the
tip of $\boldsymbol B$) has its frequency raised by $\nu_L$.

\paragraph{Dispersion.} For a damped oscillator with time dependence
$\e^{-i\omega t}$ the susceptibility near resonance is proportional to
$1/(x-ia)$ with $x = (\omega_0-\omega)/\Delta\omega_D$; averaged over a
Maxwellian, the complex refractive index of a mode is
$N-1\propto F(a,x)+iH(a,x)$: $H$ gives absorption, $F$ the phase velocity
(Kramers--Kronig). $x$ has the sign of $\lambda-\lambda_0$, so $x=v$.

\paragraph{Field along the ray ($\theta=0$).} The normal modes are the
two circular polarizations. Writing $I_\pm$ for the intensities of
positive and negative helicity, $I = I_++I_-$ and $V = I_+-I_-$, the mode
absorption gives $\eta_V = \tfrac12(H_{-1}-H_{+1})$. A linear
polarization is the sum of the two modes with equal amplitude; their
phase difference $\delta$ after a distance $s$ turns the plane of
polarization by $-\delta/2$, clockwise as seen by the observer when the
positive-helicity mode is the slower one. With the transfer equation
$\dd\vI/\dd s = -K\vI+\ldots$ and $K$ in the form of
eq.~(\ref{eq:Kmatrix}) below this gives $\rho_V = \tfrac12(F_{+1}-F_{-1})$.

\paragraph{Field across the ray ($\theta=90^\circ$, $\chi=0$).} The
modes are linear, along $\boldsymbol B$ (absorbed by the $\pi$ components,
profile $H_0$) and perpendicular to it ($\tfrac12(H_{+1}+H_{-1})$):
$\eta_Q = \tfrac12[H_0-\tfrac12(H_{+1}+H_{-1})]$. A phase lag between
them turns $U$ into $V$, and working out the sense gives
$\rho_Q = -\tfrac12[F_0-\tfrac12(F_{+1}+F_{-1})]$.

\paragraph{General direction.} With $\theta$ the angle between
$\boldsymbol B$ and the ray and $\chi$ the azimuth of its projection
from $\boldsymbol e_1$ (Stokes $Q>0$ along $\boldsymbol e_1$,
$(\boldsymbol e_1,\boldsymbol e_2,\boldsymbol n)$ right-handed, $\chi$
counterclockwise as seen by the observer), in the IAU convention
\begin{equation}
\begin{aligned}
\eta_I &= \tfrac12\bigl[H_0\sin^2\theta+\tfrac12(H_{+1}+H_{-1})(1+\cos^2\theta)\bigr],\\
\eta_{Q,U} &= \tfrac12\bigl[H_0-\tfrac12(H_{+1}+H_{-1})\bigr]\sin^2\theta\,(\cos2\chi,\sin2\chi),\\
\eta_V &= \tfrac12(H_{-1}-H_{+1})\cos\theta,\\
\rho_{Q,U} &= -\tfrac12\bigl[F_0-\tfrac12(F_{+1}+F_{-1})\bigr]\sin^2\theta\,(\cos2\chi,\sin2\chi),\\
\rho_V &= \tfrac12(F_{+1}-F_{-1})\cos\theta,
\end{aligned}
\end{equation}
all times $1/(\sqrt\pi\,\Delta v_D)$, in
\begin{equation}
K = \begin{pmatrix}
\eta_I & \eta_Q & \eta_U & \eta_V\\
\eta_Q & \eta_I & \rho_V & -\rho_U\\
\eta_U & -\rho_V & \eta_I & \rho_Q\\
\eta_V & \rho_U & -\rho_Q & \eta_I
\end{pmatrix}. \label{eq:Kmatrix}
\end{equation}

\subsection{RH's sign of Stokes V}
RH's profiles agree with these expressions once Stokes $V$ changes sign,
that is with $DKD$ in place of $K$, where $D=\operatorname{diag}(1,1,1,-1)$:
$\eta_V$, $\rho_Q$ and $\rho_U$ change sign, $\rho_V$ does not. The test finds this to rounding
($10^{-16}$) for every configuration, field and ray. RH is therefore
internally consistent, with positive $V$ meaning negative helicity:
clockwise as seen by the observer, which is ``right-handed'' in the
optics convention and what the comment in the source of \code{Zeeman()}
states: the red $\sigma$ component is right-handed for a field toward the
observer. It
is the opposite of IAU 1973. For a field toward the observer RH gives a
positive blue lobe of $V$; in the weak-field limit (below)
$V = -\bar g\,v_B\cos\theta\,\partial I/\partial v$.

For rays going down (away from the observer) RH uses the frame
$(\boldsymbol e_1',\boldsymbol e_2') = (\boldsymbol e_2,\boldsymbol e_1)$
with $\boldsymbol n' = -\boldsymbol n$, which is right-handed: $\cos\theta$
and $\cos2\chi$ change sign, $\sin2\chi$ does not. The test checks
both directions.

\subsection{Configurations}
\begin{center}
\begin{tabular}{llll}
\toprule
line & $g_l$, $g_u$ & $\bar g$ & components\\
\midrule
$^1S_0\to{}^1P_1$ (normal triplet) & --, 1 & 1 & 3\\
triplet with explicit $\bar g = 1.5$ & & 1.5 & 3\\
$^2S_{1/2}\to{}^2P_{1/2}$ (as Na D$_1$) & 2, 2/3 & 4/3 & 4\\
$^2S_{1/2}\to{}^2P_{3/2}$ (as Na D$_2$) & 2, 4/3 & 7/6 & 6\\
$^5P_2\to{}^5D_2$ & 11/6, 3/2 & 5/3 & 13 ($\pi$ $M=0\to0$ of strength 0)\\
$^5P_1\to{}^5D_0$ & 5/2, -- & 5/2 & 3\\
$^5D_0\to{}^5F_1$ (null line) & --, 0 & 0 & 3, unshifted\\
$^1S_0\to{}^3P_0$ & & & rejected\\
\bottomrule
\end{tabular}
\end{center}
The null line has $g(^5F_1) = 1+(2+6-12)/4 = 0$: no splitting, and
$Q=U=V=0$ exactly at any field strength. The $J_u=J_l=1/2$ line is not
special here: RH treats the Zeeman effect only (it is a multi-level, not
a multi-term code), and for Zeeman splitting D$_1$ is an ordinary
anomalous pattern. The field strengths (0 to 0.5\,T, $v_B$ up to
$\approx3$), inclinations (toward and away from the observer, transverse,
along the vertical) and two rays include a field exactly along the
oblique ray.

Magnetic-dipole (M1) lines, such as [O\,I] 630.03\,nm, have the same
pattern and strengths (the transition is rank 1 as well), but their
linear polarization is rotated by $90^\circ$: $\eta_{Q,U}$ and
$\rho_{Q,U}$ change sign, $\eta_V$ and $\rho_V$ do not. RH has no M1
lines (wishlist F-2); once it does, they belong in this test.

\paragraph{Defects found.} B-27: a $J=0\to0$ line was marked polarizable
($|J_u-J_l|\le1$), and \code{Zeeman()} normalized its only component, of
strength 0, by zero, making every profile of the line NaN; such lines are
now treated without polarization (\code{zeemanPolarizable()}). B-28:
\code{Bproject} obtained $\cos2\chi$ and $\sin2\chi$ by dividing by
$1-\cos^2\theta$, which is zero for a field exactly along an oblique ray:
NaN in the $Q$ and $U$ profiles of every polarizable line at that point.
The terms are multiplied by $\sin^2\theta=0$, so any finite $\chi$ will do.

\subsection{Milne--Eddington}
With RH's profiles $\Phi$ of one grid point, a slab with continuum
opacity $\chi_c$, line opacity $\chi_l\Phi$ ($\chi_l$ constant,
$\eta_0 = 5$ times the continuum at line centre) and the same source
function $S = S_0+S_1\tau_c$ for line and continuum has, in continuum
optical depth, the constant propagation matrix $\hat K = 1+\chi_l\Phi/\chi_c$
and emission vector $\hat KS\boldsymbol e_I$. By
eq.~(\ref{eq:Iplus}) in matrix form (section~\ref{sec:stokes}) the
emergent Stokes vector is the Unno--Rachkovsky solution
\begin{equation}
\vI(0) = S_0\boldsymbol e_I+\mu S_1\hat K^{-1}\boldsymbol e_I .
\end{equation}
The test builds the slab's opacity and emissivity from RH's profiles the
way \code{opacity.c} does, solves it with the polarized solver (exact for
linear $S$ and constant $K$, section~\ref{sec:stokes}) and compares across
the line with $\hat K$ from the derived profiles: $10^{-13}$ of $I$ for
all Stokes parameters.

\paragraph{Weak field.} To first order in $v_B$,
$H_{\pm1}(v)\approx H(v)\mp\bar g\,v_B\,H'(v)$ in the sense that
$\tfrac12(H_{+1}-H_{-1}) = -\bar g\,v_B H'$, so in RH's sign
$\eta_V = -\bar g\,v_B\cos\theta\,H'/(\sqrt\pi\Delta v_D)$, and
$\hat K^{-1}$ to first order gives
\begin{equation}
V = -\mu S_1\frac{\chi_l\phi_V}{(1+\chi_l\phi_I)^2}
  = -\bar g\,v_B\cos\theta\,\frac{\partial I}{\partial v},
\end{equation}
with $I = S_0+\mu S_1/(1+\chi_l\phi_I)$ the field-free intensity. In
the IAU sign this is $V = +\bar g v_B\cos\theta\,\partial I/\partial v$,
or $V = -\bar g\,\nu_L\cos\theta\,\partial I/\partial\nu$ in frequency.
At $v_B = 0.023$ RH's $V$ agrees with it to $4$--$6\times10^{-4}$ of
$\max|V|$, the size of the neglected $O(v_B^2)$ terms.

%======================================================================
\section{End to end: rhf1d with a fictional atom}\label{sec:e2e}
\code{tests/e2e/stokes\_1d.py} (\code{ctest -L e2e}) runs the whole of
\code{rhf1d}: reading the atom and the field, \code{Bproject},
\code{Zeeman}, the profiles, the opacities and the polarized solver. The
atom \code{TI\_fict.atom} has two lines from a $^2S_{1/2}$ ground level,
analogues of the sodium D lines at made-up energies (588.24 and
581.40\,nm), to $^2P_{1/2}$ ($\bar g = 4/3$) and $^2P_{3/2}$
($\bar g = 7/6$), both with anomalous patterns that RH derives from the
level labels. The atom is active with LTE populations and
\code{N\_MAX\_ITER = 0}, so RH computes one formal solution;
\code{N\_MAX\_SCATTER = 1} makes that solution use the initial $J=0$ in
the background scattering of every run. Every run then solves the same
linear problem, $S$ fixed, and the relations below hold to rounding or to
discretization error, not to the accuracy of a physical approximation.
The rays of the 1-D code lie in the $x$--$z$ plane.

\subsection{Exact symmetries}
\paragraph{No field.} $Q=U=V=0$ exactly.

\paragraph{Parity.} Reflection in the plane of the rays, $y\to-y$, maps
a solution of the transfer equation onto a solution: the Zeeman effect
and the transfer equation are invariant under parity. $\boldsymbol B$ is
an axial vector, so the reflected field is
$-(B_x,-B_y,B_z) = (-B_x,B_y,-B_z)$, i.e.\ $\gamma\to\pi-\gamma$,
$\chi\to\pi-\chi$. The reflection reverses the helicity ($V\to-V$) and
mirrors the direction of linear polarization about $\boldsymbol e_1$,
which lies in the plane of the rays ($U\to-U$). The test compares
$(I,Q,-U,-V)$ of the reflected field with the original on FAL-C with a
field that varies with depth (200\,G to 2\,kG, $\gamma$ from 70 to
$20^\circ$, $\chi$ from 100 to $10^\circ$): $3\times10^{-13}$, with and
without magneto-optical effects.

\paragraph{Reversed field.} $\boldsymbol B\to-\boldsymbol B$
($\gamma\to\pi-\gamma$, $\chi\to\chi+\pi$) changes the sign of $\eta_V$
and $\rho_V$ and leaves the other elements of $K$. Without magneto-optical
terms this is $K\to DKD$ with $D=\operatorname{diag}(1,1,1,-1)$, so the
solution is $(I,Q,U,-V)$: $2\times10^{-13}$. With them it is no symmetry:
$DKD$ would also change $\rho_Q$ and $\rho_U$ and keep $\rho_V$, and
Faraday rotation, odd in $B$, turns the plane of linear polarization the
other way (on FAL-C $Q$ changes by 66\% of its maximum).

\subsection{Turning the reference frame}
For a vertical ray, turning the field azimuth by $\Delta$ turns the frame:
$K\to RKR^{\mathsf T}$ with $R$ the rotation by $2\Delta$ in the $(Q,U)$
plane, and the transfer equation gives $(Q,U)\to R(Q,U)$ with $I$ and $V$
unchanged. Steffen's limiter (\code{cent\_deriv}) is nonlinear, so
applied to each element of $\vS_K'$ and $\dot K$ on its own it does not
commute with mixing the $Q$ and $U$ components (N-25). Take $Q$ with a
maximum at a point (one-sided slopes $1$ and $-1$: derivative $0$) and $U$
rising (slopes $1$ and $5$: derivative $2$). Limited first and then turned
by $45^\circ$ in the $(Q,U)$ plane (the frame by $22.5^\circ$, so that
$Q'=(Q+U)/\sqrt2$), $Q'$ gets the derivative $(0+2)/\sqrt2\approx1.41$; turned
first, $Q'$ has slopes $2/\sqrt2$ and $4/\sqrt2$ and gets $3/\sqrt2\approx2.12$.
The solver therefore limits the pairs that turn with the frame,
$(S_Q,S_U)$, $(\eta_Q,\eta_U)$ and $(\rho_Q,\rho_U)$ (the elements
$K_{01},K_{02}$; $K_{10},K_{20}$; $K_{13},K_{23}$; $K_{31},K_{32}$), as
vectors. With one-sided slopes $\mathbf a$, $\mathbf b$ and the weighted
centred slope $\mathbf p$,
\begin{equation}
\mathbf d = \begin{cases}
0, & \mathbf a\cdot\mathbf b \le 0,\\
\mathbf p\,\min\bigl(1,\ 2|\mathbf a|/|\mathbf p|,\ 2|\mathbf b|/|\mathbf p|\bigr), & \text{otherwise},
\end{cases}
\end{equation}
which uses only lengths and a dot product, so it turns with the frame, and
which for a single component is Steffen's limiter. On FAL-C (82 points)
with the depth-dependent field and $\Delta = 35^\circ$, $Q$ and $U$ now
equal the rotated ones to $3\times10^{-14}$ of their maximum, $V$ to
$2\times10^{-14}$, $I$ to $2\times10^{-15}$; on the grey atmosphere below
to $10^{-14}$ at every $N$. Limited element by element they differed by
$7\times10^{-3}$, $2\times10^{-4}$ and $10^{-5}$ on FAL-C, and on the grey
atmosphere by $4.8\times10^{-3}$ ($N=40$) to $5.0\times10^{-6}$ ($N=320$),
a discretization error close to third order. The test requires the
rotation to hold to $10^{-12}$ on both.

\subsection{Weak field without a wavelength derivative}
Measure wavelength in velocity units, $u = c(\lambda-\lambda_0)/\lambda_0$.
To first order in $B$, for a field constant along the ray, the
transfer equation for $V$ is
\begin{equation}
\frac{\dd V}{\dd s} = -\eta_I V - \eta_V(I^{(0)}-S),
\end{equation}
with $I^{(0)}$ the field-free intensity and, in RH's sign
(section~\ref{sec:zeeman}),
$\eta_V=-\bar g\,u_B\cos\theta\,\partial\eta_I/\partial u$, where
$u_B = e\lambda_0B/(4\pi m_e)$ is the Zeeman splitting in velocity, the
same at every depth. Shift the line profiles rigidly by a velocity
$\sigma$, the same at every depth; then $\partial I/\partial\sigma$ obeys
\begin{equation}
\frac{\dd}{\dd s}\frac{\partial I}{\partial\sigma}
= -\eta_I\frac{\partial I}{\partial\sigma}
 - \frac{\partial\eta_I}{\partial\sigma}(I^{(0)}-S),
\qquad \frac{\partial\eta_I}{\partial\sigma} = -\frac{\partial\eta_I}{\partial u},
\end{equation}
with the same boundary condition ($S$ does not depend on $\sigma$: LTE,
and $J=0$ in the background scattering). The two equations are the same,
so
\begin{equation}
V = \bar g\,u_B\cos\theta\,\frac{\partial I}{\partial\sigma}
\end{equation}
for any depth dependence of $T$, $n_e$, the Doppler width and the
opacities; where the continuum does not vary across the line,
$\partial I/\partial\sigma = -\partial I/\partial u$ and this is the
familiar $V=-\bar g\,u_B\cos\theta\,\partial I/\partial u$. A uniform
vertical velocity $v_z$ is such a shift: for an outgoing ray at $\mu$,
$\sigma=-\mu v_z$ (to the blue for an upflow; \code{profile.c} evaluates
the profile at $v+\mu v_z/v_D$ in Doppler units), so
\begin{equation}
V = -\bar g\,u_B\,\frac{\cos\theta}{\mu}\,\frac{\partial I}{\partial v_z}.
\end{equation}
The test takes $\partial I/\partial v_z$ from four field-free runs with
$v_z=\pm20$ and $\pm40\,\mathrm{m\,s^{-1}}$ (fourth-order central
difference, error $\sim(\delta v/v_D)^4\approx10^{-8}$), and $V$ from a run
with $B=1$\,G, $\gamma=30^\circ$, $\chi=20^\circ$, where the neglected
$O(B^2)$ terms are $(u_B/v_D)^2\sim10^{-7}$. No wavelength derivative is
taken, and nothing in the comparison depends on the wavelength grid.

The two sides reach the solver by different routes, an extra source term
in the polarized system against shifted opacities, so they agree to
discretization error. The test uses a grey atmosphere with no kinks,
$T^4=\tfrac34T_\text{eff}^4(\tau+\tfrac23)$ with $\tau$ the column mass in
g\,cm$^{-2}$, $\log m$ from $-5$ to 1 at $N$ points, $n_e$ from
\code{SOLVE\_NE}. For the three rays and both lines the largest
$|V-V_\text{wf}|/\max|V|$ is
\begin{center}
\begin{tabular}{rcccc}
\toprule
$N$ & 40 & 80 & 160 & 320\\
\midrule
error & $5.6\times10^{-3}$ & $4.2\times10^{-4}$ & $2.7\times10^{-5}$ &
$1.9\times10^{-6}$\\
\bottomrule
\end{tabular}
\end{center}
fourth order, as for the 1-D Bezier solver itself
(section~\ref{sec:bezier}). This checks, end to end, the sign of $V$,
$\bar g$ from the labels, the Larmor splitting, $\cos\theta$ for oblique
rays and the units of the field input.

%======================================================================
\section*{Regenerating the reference values}
In \code{tests/unit}, each script writes a header to standard output:
\begin{itemize}
\item \code{make\_voigt\_reference.py} $\to$ \code{voigt\_reference.h}
\item \code{make\_kernels\_reference.py} $\to$ \code{kernels\_reference.h}
\item \code{make\_bezier\_reference.py} $\to$ \code{bezier\_reference.h}
\item \code{make\_weideman\_coefficients.py} $\to$ the coefficient table
  in \code{voigt.c}
\end{itemize}
All need \code{mpmath}.

\bibliographystyle{plainnat}
\bibliography{unit-tests}

\end{document}
